\documentclass{article}

\usepackage[utf8]{inputenc}
\usepackage{amsmath, esint}
\usepackage{wasysym}
\usepackage{qrcode}
\usepackage[colorlinks]{hyperref}
\usepackage{lmodern}
\usepackage{graphicx}
\usepackage{xcolor}
\usepackage[left=2cm, top=3cm, right=2cm]{geometry}
\usepackage{minted}
\usepackage{booktabs}
\usepackage{svg}
\usepackage{xcolor}
\definecolor{LightGray}{gray}{0.975}
\hypersetup{
  urlcolor=blue
}

\title{ML101 \\ Homework 01: A `gentle' Introduction to Machine Learning.}
\author{Andrés Oswaldo Calderón Romero, Ph.D.}
\date{\today}

\begin{document}

\maketitle

\section{Introduction}
\label{sec:intro}

Machine Learning relies on a blend of statistical foundations, algorithmic design, and domain intuition. To kick off the semester, this assignment provides a structured framework to explore your first textbook chapter.

Through this task, you will engage with state-of-the-art AI tools to extract key takeaways, probe deeper into complex concepts through targeted questions, and visually organize your insights in a handwritten mind map. This approach is intended to strengthen both your technical curiosity and your ability to critically evaluate AI-generated knowledge.


\section{Reading Material}
\label{sec:reading}

We will select \textbf{one} of the following textbooks to cover the foundations and core principles of machine learning:

\begin{enumerate}
    \item \textit{Understanding Machine Learning: From Theory to Algorithms} by Shai Shalev-Shwartz and Shai Ben-David.
    \item \textit{An Introduction to Statistical Learning: with Applications in R} (or Python) by Gareth James, Daniela Witten, Trevor Hastie, Robert Tibshirani, and Jonathan Taylor.
    \item \textit{Probabilistic Machine Learning: An Introduction} by Kevin P. Murphy.
\end{enumerate}

\section{Key Takeaways}
\label{sec:keypoints}

We will use an AI tool to summarize Chapter 1 of our selected textbook. Feel free to request a translation of the content if you prefer reading in another language. Please keep a record of all prompts used to generate and refine your summary.

Additionally, ask the AI three (3) follow-up questions to explore specific concepts in greater depth. Be sure to document the exact prompts used for each question.

\section{Individual Assignment}
\label{sec:work}

Based on the summary and the answers to your follow-up questions, create a mind map to illustrate your key findings and main concepts. The mind map must be \textbf{handwritten}.

\section{Expected Deliverable}
\label{sec:deliverables}

You will prepare a well-structured report documenting and discussing all prompts used in Section~\ref{sec:keypoints}, along with a photo or scan of your handwritten mind map from Section~\ref{sec:work}.

Please submit your report in \textbf{PDF format} through the corresponding Brightspace\texttrademark{} assignment portal by \textbf{August 13, 2026}.

\vspace{5mm}
Happy Reading! \includesvg[width=4mm]{figures/sunglasses}

\end{document}

